% Generated by tools/edn_latex.py from reports/edn.md.
% Do not edit: change reports/edn.md and run `make edn`.
\ednentry{
  id = {EDN-74},
  title = {A run is tagged with the hashes of its own stage's inputs, and \texorpdfstring{\texttt{git\_\allowbreak{}dirty}}{git\_dirty} ignores what DVC writes itself},
  date = {2026-10-05},
  milestone = {M3: Quality Assurance},
  topic = {Experiment Tracking, Reproducibility},
  participants = {@lukas2510, decided by the repository owner},
  decision = {Every MLflow run a pipeline stage opens or appends to is tagged with the hash of each dependency \texttt{dvc.\allowbreak{}yaml} declares for that stage, one \texttt{<stage>.\allowbreak{}deps.\allowbreak{}<path>} tag each, computed through DVC's own API when the stage starts; it is the value \texttt{dvc.\allowbreak{}lock} records under \texttt{stages.\allowbreak{}<stage>.\allowbreak{}deps} once the stage has finished.
\texttt{git\_\allowbreak{}dirty} is computed while ignoring the files \texttt{dvc repro} writes itself: each pipeline's lock file and every output declared \texttt{cache: false}, read off the pipeline rather than listed.
\texttt{dvc\_\allowbreak{}lock\_\allowbreak{}md5}, the MD5 of the whole lock, is dropped.
\texttt{evaluate}, which appends to the run \texttt{train} opened (EDN-55), records its own \texttt{evaluate.\allowbreak{}git\_\allowbreak{}commit}, \texttt{evaluate.\allowbreak{}git\_\allowbreak{}dirty} and \texttt{evaluate.\allowbreak{}deps.\allowbreak{}<path>} beside the fit's \texttt{git\_\allowbreak{}commit}, \texttt{git\_\allowbreak{}dirty} and \texttt{train.\allowbreak{}deps.\allowbreak{}<path>}.
Input hashes are recorded only when DVC runs the stage, which DVC signals by setting \texttt{DVC\_\allowbreak{}STAGE}.},
  rationale = {\emph{Alternatives considered:}
\begin{itemize}[leftmargin=*, topsep=2pt]
\item \textbf{Option A (chosen): the hashes of the stage's own inputs, taken when it starts, and \texttt{git\_\allowbreak{}dirty} without DVC's writes.}
\newline Pros: every tag names something a commit or the DVC remote reproduces: for the runs of a committed \texttt{dvc repro}, each input tag equals the committed lock line for line, and a clean commit is reported clean.
The tags exist from the moment the run does, so a stage that fails half-way still says what it read.
A data dependency's hash is its address in the DVC cache and on the remote, so \texttt{dvc pull} restores that input from it and it identifies the input whatever happens to the commits, a squash merge or a rebase included; a code file is restored by git at the commit, where a clean tree guarantees it has the recorded hash.
And one tag per input lets MLflow find every run that read a given matrix with an exact filter.
\newline Cons: a run carries 24 provenance tags instead of three, ten input hashes and a commit and a dirty flag for each of its two stages; the stage opens the DVC repository in-process, which costs about 0.3 s per run, measured at 0.28 to 0.41 s on the laptop against fits of 0.1 to 19 s; the code depends on DVC's Python API (\texttt{Index.\allowbreak{}stages}, \texttt{Output.\allowbreak{}get\_\allowbreak{}hash}), which DVC does not promise to keep stable, so a DVC upgrade can break it -- the test that runs a real \texttt{dvc repro} and compares with the lock is what would say so; and a hand edit of an uncached output, \texttt{metrics.\allowbreak{}json} say, no longer marks a run dirty, which costs nothing today, because no stage that opens or appends to a run depends on an uncached output; a stage that did would hash it among its inputs, so the edit would show there.
\item \textbf{Option B: keep the whole-lock MD5, and document what it does not identify.}
\newline Pros: no code change; the model card already warned that the four runs of one ladder carry four different digests.
\newline Cons: the tag identifies no committed state, so documenting it would only document that NFR-14's data-version clause is unmet.
Measured on the runs \texttt{main}'s models pointed at before this change: four \texttt{train} runs, four different digests, none equal to the committed lock; and the first stage's digest equals the lock from before the re-run.
It also leaves \texttt{git\_\allowbreak{}dirty} true for every stage after the first of a clean \texttt{dvc repro}, three of those four runs, only because \texttt{dvc.\allowbreak{}lock} had changed under them, so the report could not say its runs were made from a committed state either.
\item \textbf{Option C: tag after the run instead of before it.}
Either at the end of the stage's own process, or by a separate step after \texttt{dvc repro} that reads the committed lock and writes the tags onto the runs.
\newline Pros: the separate step reads the very lock it is compared with, so the tags would equal it by construction; and the stage-end variant needs no change to how the run is opened.
\newline Cons: at the end of its own process a stage still cannot read its lock entry, because DVC writes it only after the command has exited, so the stage-end variant would have to compute the same hashes A computes, later and with less of the run covered: a stage that fails half-way would carry none.
The separate step is a second command somebody has to remember after every \texttt{dvc repro}, and one forgotten leaves the runs without provenance and nothing failing; it would also describe the tree when the step ran, not when the stage did, and attribute anything changed in between to the run.
And it equals the lock by copying it, so it could not detect a run whose inputs differ from what the lock says, which is the one thing the check is for.
\end{itemize}
\par
\emph{Why:} NFR-14 promises that a run can be traced back to what produced it, and only a committed state can be traced back to.
A whole-lock digest taken during a \texttt{dvc repro} is never the lock committed with that run, so option B keeps a tag that cannot do its job, and C either recomputes A's hashes later or copies the lock instead of checking against it.
A is the only option whose tags are evidence: two programs, the stage at its start and DVC at its end, write the same values independently, and the check is that they agree.
\par
Three implementation choices are part of the decision.
The hashes are computed through DVC's API rather than read from upstream stages' outputs in \texttt{dvc.\allowbreak{}lock}, because only some dependencies are another stage's output: the code files are no stage's output, and \texttt{evaluate}'s \texttt{models} and \texttt{data/\allowbreak{}processed/\allowbreak{}features} are the parents of several outputs and the output of none, so their \texttt{.\allowbreak{}dir} hashes are in the lock only as \texttt{evaluate}'s own dependencies, written after it has run.
Reimplementing DVC's directory hash with \texttt{hashlib} would be one more thing that can disagree with the lock; calling the code that writes the lock cannot, and \texttt{tests/\allowbreak{}test\_\allowbreak{}provenance.\allowbreak{}py} proves the two equal on a three-stage pipeline DVC really runs, on a directory, a file, a \texttt{foreach} stage and an uncached output that a later stage depends on.
The tags are recorded only under DVC because every test calls the stages with paths under a temporary directory: the hashes of the files \texttt{dvc.\allowbreak{}yaml} declares would then describe files the run never opened.
And \texttt{evaluate} records its own commit beside the fit's because the two need not run at the same commit: a change to \texttt{evaluate.\allowbreak{}py} alone reruns only \texttt{evaluate}, against the runs of an earlier \texttt{train}, and overwriting \texttt{git\_\allowbreak{}commit} would then claim the fit came from code it did not come from.
\par
The implementation found more than the issue described.
DVC deletes a stage's outputs before running it, so a stage with an uncached output starts with that git-tracked file deleted: on a three-stage pipeline reproduced from a clean commit with the old code, even the first stage was tagged dirty, and only the first stage's digest matched any committed lock, the one from before the run.
And DVC adds a line to the \texttt{.\allowbreak{}gitignore} beside a cached output the first time it produces it, which for \texttt{gx/\allowbreak{}} is the repository root's hand-written \texttt{.\allowbreak{}gitignore}; ignoring that file would hide every edit to it, so the first run of a new cached output is reported dirty until its line is committed, which \texttt{pipeline.\allowbreak{}md} says.
And DVC addresses a stage relative to the directory \texttt{dvc repro} was started in, \texttt{.\allowbreak{}.\allowbreak{}/\allowbreak{}dvc.\allowbreak{}yaml:train@\allowbreak{}b0} from a subdirectory, so the stage is looked up by its name alone; an independent review of the pull request found that the first version resolved the path against the repository root, and that from a subdirectory every run would then have lost its tracking, under a warning that blamed the credentials.
\par
Checked on the real snapshot rather than only in a test: the four \texttt{train} stages and \texttt{evaluate} were re-run from a clean commit with \texttt{RECOMMENDITOS\_\allowbreak{}REQUIRE\_\allowbreak{}TRACKING=\allowbreak{}1}, and \texttt{reports/\allowbreak{}analysis/\allowbreak{}run\_\allowbreak{}provenance\_\allowbreak{}check.\allowbreak{}py} read the four runs back from DagsHub.
Before the re-run, the four runs \texttt{main}'s committed models pointed at carried none of the input tags, three of them were tagged dirty, and each \texttt{dvc\_\allowbreak{}lock\_\allowbreak{}md5} differed from the digest of the committed lock.
After it, each of the four carries ten \texttt{train.\allowbreak{}deps} and ten \texttt{evaluate.\allowbreak{}deps} tags, every one equal to the committed lock, and \texttt{git\_\allowbreak{}dirty} and \texttt{evaluate.\allowbreak{}git\_\allowbreak{}dirty} are \texttt{false} for all four, naming the clean commit the re-run started from (\texttt{reports/\allowbreak{}analysis/\allowbreak{}run\_\allowbreak{}provenance\_\allowbreak{}check\_\allowbreak{}results.\allowbreak{}txt}).
The re-run, 1 min 48 s at a peak RSS of 693 MB, moved no number: \texttt{metrics.\allowbreak{}json}, the segment table and the masking sweep are byte-identical, each per-variant record changed only in its MLflow run id, and inside each model bundle only \texttt{model.\allowbreak{}json} changed, by the provenance EDN-56 puts there, while every payload file kept its hash.},
  ai = {Information seeking, Alternative generation, Alternative assessment, Recommendation, Solution generation},
  response = {Accepted with modifications},
  assessment = {AI found the defect in an independent review of PR \#70, by reading the runs back from DagsHub instead of trusting the tests, which checked that the tags exist and that they change when \texttt{dvc.\allowbreak{}lock} changes, not that they identify a committed state.
It measured the four digests against the committed lock and the dirty flags against the files that had changed, generated options A to C and assessed them, and recommended A, fixed now in a pull request of its own; Lukas chose that, before the first delivery, so that the report's provenance claim is true and the CodeCarbon run of \#38 retrains only once afterwards.
The implementation details were AI's, as were the two further findings, the deleted uncached output and the \texttt{.\allowbreak{}gitignore} line, both reproduced before they were acted on.
An independent AI review of the pull request reproduced the evidence and found the subdirectory regression, two overstated claims -- that every hash retrieves its input from the remote, and that a hand edit of the validation summary would show in a run -- and three smaller gaps: a dependency path MLflow refuses as a tag key, a failed \texttt{git status} read as a clean tree, and a test pipeline without a parameter; each was fixed, test first where there was code, before merging, which is the modification.},
  aievidence = {Claude Code sessions on 2026-10-05: the independent review of PR \#70, whose measurement is the body of issue \#79, and the implementation of issue \#79, in which the issue's scenario was reproduced on a three-stage pipeline with \texttt{main}'s \texttt{run\_\allowbreak{}provenance} (every stage tagged dirty, the digests of the lock before the run and of two mid-run locks), the new tests were confirmed to fail on \texttt{main}, where the module they test does not exist, and the check script was run against DagsHub before and after the re-run.},
  evidence = {\href{https://github.com/mlops-2627q1-mds-upc/MLOPS_Recommenditos/issues/79}{issue \#79}; \href{https://github.com/mlops-2627q1-mds-upc/MLOPS_Recommenditos/pull/80}{PR \#80}; \href{https://github.com/mlops-2627q1-mds-upc/MLOPS_Recommenditos/blob/main/recommenditos/provenance.py}{\texttt{recommenditos/\allowbreak{}provenance.\allowbreak{}py}} and \texttt{recommenditos/\allowbreak{}tracking.\allowbreak{}py}; \href{https://github.com/mlops-2627q1-mds-upc/MLOPS_Recommenditos/blob/main/tests/test_provenance.py}{\texttt{tests/\allowbreak{}test\_\allowbreak{}provenance.\allowbreak{}py}} and the NFR-14 tests of \texttt{tests/\allowbreak{}test\_\allowbreak{}tracking.\allowbreak{}py}; \href{https://github.com/mlops-2627q1-mds-upc/MLOPS_Recommenditos/blob/main/reports/analysis/run_provenance_check.py}{\texttt{reports/\allowbreak{}analysis/\allowbreak{}run\_\allowbreak{}provenance\_\allowbreak{}check.\allowbreak{}py}} and \href{https://github.com/mlops-2627q1-mds-upc/MLOPS_Recommenditos/blob/main/reports/analysis/run_provenance_check_results.txt}{its results}; \href{https://github.com/mlops-2627q1-mds-upc/MLOPS_Recommenditos/blob/main/docs/docs/pipeline.md}{The DVC pipeline}, From a run to its inputs; \href{https://github.com/mlops-2627q1-mds-upc/MLOPS_Recommenditos/blob/main/docs/docs/specification.md}{specification} NFR-14; \hyperref[EDN-55]{EDN-55}, \hyperref[EDN-66]{EDN-66}.},
}
